\section{结论}\label{conclusion}

计算机的主要性能指标包括可解决问题的规模和范围、求解的速度和精度，以及计算机的物理尺寸、可靠性和成本. 这些指标之间存在强烈的相互制约关系，不太可能有一种计算机形式在所有指标上都达到最优. 复杂过程的辨识与仿真，以及多变量控制系统的实现，需要大量的计算元件（如乘法器、加法器和积分器）同时工作且成本低廉. 然而，这些元件并不需要快速或精确地计算解，因在经济和化学过程的仿真中，10 cs 的带宽和 1\% 的总体精度已经足够，而在存在反馈的控制系统中，10\% 的计算精度可能就已充足. 在这些情况下，用数字计算机的精度和模拟计算机的速度来换取随机计算机的经济性是较为有利的.

\section{致谢}\label{acknowledgments}

感谢 J. H. Andreae 博士（现就职于新西兰基督城坎特伯雷大学）就学习机和计算机问题与我进行的长久讨论，以及对本文稿的批评意见；感谢 P. L. Joyce 先生构建了第一台随机计算机并获得了图 9 的实验结果；感谢 S.T.L. 允许发表本文.

\section{参考文献}\label{references}

1 J. H. ANDREAE Learning machines in Encyclopaedia of Information,
Linguistics and Control (Pergamon Press, to be published)\\
2 B. R. GAINES and J. H. ANDREAE A learning machine in the context of
the general control problem Proceedings of the 3rd Congress of IFAC
1966\\
3 G. NAGY A survey of analog memory devices IEEE Trans. Electron. Comp.
12 388 1963\\
4 B. R. GAINES Stochastic computers in Encyclopaedia of Information,
Linguistics and Control Pergamon Press, to be published\\
5 B. R. GAINES Techniques of identification with the stochastic computer
Proc. IFAC Symp. Problems of Identification 1967\\
6 F. V. MAYOROV and Y. CHU Digital differential analysers IIiffe Books,
London 1964\\
7 H. SCHMID ``An operational hybrid computing system IEEE Trans.
Electron Comp. 12 715 1963\\
8 B. R. GAINES and P. L. JOYCE Phase computers 5th AICA Congress 1967\\
9 W. W. PETERSON Error correcting codes MIT Press \& Wiley, New York
1961\\
10 C. L. BECKER and J. V. WAIT Two-level correlation on an analog
computer IRE Trans. Electron Comp. 10 752 1961

11 B. P. TH. VELTMAN and A. van den BOS The applicability of the relay
correlator and the polarity coincidence correlator in automatic control
Proc. 2nd Congress IFAC 1963\\
12 P. EYKHOFF, P. M. van der GRINTEN, H. KWAKERNAAK, B. P. TH. VELTMAN
Systems modelling and identification Survey Paper 3rd Congress of IFAC
1966\\
13 O. I. ELGERD High frequency signal injection: a means of changing the
transfer characteristics of nonlinear elements WESCON 1962\\
14 A. A. PERVOZANSKII Random processes in nonlinear control Academic
Press, New York 1965\\
15 P. JESPERS, P. T. CHU, and A. FETTWEIS A new method to compute
correlation functions Proc. Int. Symp. Inf. Theo. 1962\\
16 G. A. KORN Random-process simulation and measurement McGraw Hill,
Inc.~1966\\
17 M. O. RABIN Probabilistic automata Inf. \& Contr. 6 230 1963\\
18 A. PAZ Some aspects of probabilistic automata Inf. \& Contr. 9 26
1966\\
19 F. ROSENBLATT A model for experimental storage in neural networks in
Computer and Information Sciences Spartan Books, Washington, D.C. 1964\\
20 B. WIDROW and F. W. SMITH Pattern-recognizing control systems in
Computer and Information Sciences Spartan Books, Washington, D.C. 1964\\
21 A. NOVIKOFF Convergence proofs for perceptrons in Mathematical Theory
of Automata Polytechnic Press, Brooklyn \& Wiley Interscience 1963
